\begin{figure}[ht]
\centering
\begin{tikzpicture}[x=1.2cm, y=1.0cm, font=\small]

% Axes
\draw[->, thick] (-0.3, 0) -- (5.4, 0) node[right, font=\scriptsize] {$x$};
\draw[->, thick] (0, -0.3) -- (0, 4.6) node[above, font=\scriptsize] {$f(x)$};

% Convex curve f(x) = 0.28 (x-1)^2 + 0.4, drawn over [0.4, 4.6].
\draw[engblue, very thick, smooth, samples=60, domain=0.4:4.6]
  plot ({\x}, {0.28*(\x - 1)*(\x - 1) + 0.4});
\node[engblue, anchor=south east, font=\scriptsize] at (4.55, 3.5)
  {convex $f$};

% Two sample points x_a = 1, x_b = 4.
% f(1) = 0.4, f(4) = 0.28*9 + 0.4 = 2.92.
\fill[engred] (1.0, 0.40) circle (1.6pt);
\fill[engred] (4.0, 2.92) circle (1.6pt);

% Chord between (1, 0.4) and (4, 2.92): the secant line.
\draw[engred, thick] (1.0, 0.40) -- (4.0, 2.92);
\node[engred, anchor=south, font=\scriptsize] at (2.5, 1.85) {chord};

% E[X] taken at x = 2.5 (midpoint of the two-point distribution).
% Point on the curve: f(2.5) = 0.28*2.25 + 0.4 = 1.03  -> f(E[X])
% Point on the chord at x = 2.5: (0.4 + 2.92)/2 = 1.66  -> E[f(X)]
\draw[dashed, enggray] (2.5, 0) -- (2.5, 1.66);
\fill[engorange] (2.5, 1.03) circle (1.8pt);
\fill[engorange] (2.5, 1.66) circle (1.8pt);

% Bracket showing the Jensen gap between the two heights.
\draw[decorate, decoration={brace, amplitude=4pt}, enggray]
  (2.62, 1.03) -- (2.62, 1.66)
  node[midway, anchor=west, font=\scriptsize, xshift=3pt]
  {Jensen gap};

\node[anchor=north, font=\scriptsize] at (2.5, -0.08) {$\E[X]$};
\node[engorange, anchor=east, font=\scriptsize] at (2.40, 1.03)
  {$f(\E[X])$};
\node[engorange, anchor=north west, font=\scriptsize] at (2.62, 2.15)
  {$\E[f(X)]$};

\node[engred, anchor=north, font=\scriptsize] at (1.0, -0.08) {$x_{a}$};
\node[engred, anchor=north, font=\scriptsize] at (4.0, -0.08) {$x_{b}$};

\end{tikzpicture}
\caption[Jensen's inequality for a convex function]{Jensen's
inequality read off a two-point distribution. A random variable $X$
takes the value $x_{a}$ or $x_{b}$ with equal probability, so its mean
$\E[X]$ sits at the midpoint. The value $f(\E[X])$ lies on the curve
(lower orange dot); the mean $\E[f(X)]$ lies on the chord at the same
abscissa (upper orange dot), because averaging the two function values
$f(x_{a})$ and $f(x_{b})$ is exactly evaluating the chord at the mean.
Convexity puts the chord above the curve, so $f(\E[X]) \leq \E[f(X)]$,
and the vertical bracket is the gap. The same picture drives the
arithmetic-mean / geometric-mean bound (take $f = -\log$), the
nonnegativity of variance (take $f(x) = x^{2}$), and the
evidence lower bound of variational inference (take $f = \log$, a
concave function, which flips the inequality).}
\label{fig:vol02:ch15:jensen-inequality}
\end{figure}
